\section{Espaços Conexos}
 
\begin{definition}
    Um espaço \((X,\tau)\) é dito \emph{desconexo} se \(\exists A\ab, B\ab \in \tau \setminus \{\varnothing\} : X = A \cup B\) e \(A \cap B = \varnothing\), nesse casso \(X = A \mid B\) é uma \emph{separação} de \(X\). 
\end{definition}

\begin{proposition}
    \textbf{26.3.} Cada \([a,b] \subset (\R, \tau_2) : a < b\) é conexo. \comentario{Suponha que \(\exists A \mid B = [a,b], \ b \in B\), \(c := \sup A\) e consiga a contradição}
\end{proposition}

\begin{proposition}
    \textbf{26.4.} \(X\) é conexo \(\ \leftrightharpoons \ \) \(A = \overline{A} = A^\circ \subset X \) implica \(A = \varnothing\) ou \(A = X\). \comentario{Mogollas}
\end{proposition}

\begin{proposition}
    \textbf{26.5.} Imagem contínua de conexo é conexo. \comentario{Mogollas}
\end{proposition}

\hypertarget{prop266}{}
\begin{proposition}
    \textbf{26.6.} Se \(S \subset X\)  é conexo, então \(\overline{S}\) também é conexo. \comentario{Supondo que \(\exists A \mid B = \overline{S}\), então \((S \cap A) \mid (S\cap B) = S\)\Lightning}
\end{proposition}

\hypertarget{coro267}{}
\begin{corollary}
    \textbf{26.7.} Se \(S \subset X\) é conexo e \(S \subset T \subset \overline{S}\), então \(T\) é conexo. \comentario{\(X = T \) em \(\square\) \hyperlink{prop266}{26.6.}}
\end{corollary}

%\begin{definition}
%    \(A, B \subset (X, \tau)\) são \emph{mutuamente separados} se \(\overline{A} \cap B = A \cap \overline{B} = \varnothing\). 
%\end{definition}

%\begin{proposition}
%    \textbf{26.9.} \(X\) é disconexo  \(\ \leftrightharpoons \ \) \(\exists A, B\subset X\) mutuamente separados \(: X = A \cup B\). 
%\end{proposition}

%\demo{(\(\Rightarrow\)) \(\exists A \mid B = X\), com \(A\) e \(B\) abertos e fechados, em particular, mutuamente separados; (\(\Leftarrow\)) \(\exists A, B \subset X\) mutuamente separados \(: X = A\cup B = \overline{A} \cup B = A \cup \overline{B}\), segue que \(A\mid B = X\). }

\begin{lemma}
    \textbf{*26.11.} Se \( X = A \mid B\) e \(S \subset X\) é conexo, então \(S \subset A\) ou \(S \subset B\). \comentario{Se não \(S = S \cap X = (S \cap A)\ab \cup (S \cap B)\ab \ \leftrightharpoons \ S\cap A = \varnothing\) ou \(S\cap B = \varnothing\)}
\end{lemma}

\hypertarget{prop2612}{}
\begin{proposition}
    \textbf{26.12.} Sejam \((S_i)_{i\in I} : X \supset S_i \) é conexo \(\forall i \in I,  \ X = \bigcup S_i\) e \(\bigcap S_i \neq \varnothing\), então o \(X\) todo é conexo. \comentario{Supõe que \(\exists A \mid B = X\) e seja \(s \in \bigcap S_i\) então \(S_i \subset A\) ou \(S_i \subset B, \ \forall i \in I\) logo \(A = \varnothing \) ou \(B = \varnothing\)} 
\end{proposition}

\begin{proposition}
    \textbf{26.13.} Se \(\forall x,y \in X, \ \exists S_{xy} \subset X \) conexo \(: x,y \in S_{xy}\), então \(X \) é conexo.  \comentario{Fixo \(a \in X, \ X= \underset{x \in X}{\bigcup} S_{ax}\) e \(a \in \underset{x\in X}{\bigcap}S_{ax}\), logo \(X\) é conexo, \(\square\) \hyperlink{prop2612}{26.12.} }
\end{proposition}

\begin{proposition}
    \textbf{26.14.} Se \(X = \displaystyle\bigcup_{n\in \N} S_n : S_n \cap S_{n+1} \neq \varnothing, \ \forall n \in \N\), então \(X\) é conexo. %\comentario{}
\end{proposition}

\demo{Por indução e \(\square\) \hyperlink{prop2612}{26.12.} que \(T_n := \overset{n}{\bigcup}\  S_k\) é conexo. Logo, \(X = \bigcup T_n\) e \(\bigcap T_n = T_1 \neq \varnothing\), segue mais uma vez da mesma proposição que \(X\) é conexo.  }

\begin{example}
    \(\R = \displaystyle\bigcup_{n\in \N}[-n,n]\) e \(\R^n = \displaystyle\bigcup \{\ell : \ell \text{ é reta e \textbf{0}} \in \ell\}\) são conexos.   
\end{example}

\begin{proposition}
    \textbf{26.16.} \(X:= \prod X_i\) é conexo \(\ \leftrightharpoons \ \) \(X_i \) é conexo, \(\forall i \in I\). 
\end{proposition}

\newcommand{\p}{\textbf{a}}
\demo{(\(\Rightarrow\)) \(\operatorname{Im} \pi_i = X_i\) conexo; (\(\Leftarrow\)) Demonstração da profe (mais bonita): (i) Mostra que \(X_1 \times X_2 \) é conexo; (ii) Fixo \(\p \in X, \ \forall J\subset I\) finito seja 
\[X_J := \prod_{i \in I\setminus J}\{a_i\}\times \prod_{j \in J} X_j \simeq \prod_{j\in J} X_j \ \text{ conexo \ (indução)}.\]
\(Y:= \bigcup X_J\) é conexo pois \(\p \in \bigcap X_J\), \(\square\) \hyperlink{prop2612}{26.12.}; (iii) Mostra que \(\overline{Y} = X\). 
}

\subsection*{Componentes conexas}

\begin{definition}
    Dado \(x \in (X, \tau)\) denotamos por \(C_x := \bigcup \{S \subset X : x \in S \text{ conexo}\}\), o maior subconjunto conexo que contém \(x\) (\emph{componente conexa} de \(x\)). 
\end{definition}

\begin{proposition}
    \textbf{27.2.} Dados \(x,y \in X\) tem-se que \(C_x = C_y\) ou \(C_x \cap C_y = \varnothing\). \comentario{\(\square\) \hyperlink{prop2612}{26.12.} } 
\end{proposition}

\begin{proposition}
    \textbf{27.3.} As componentes conexas são sempre fechadas. \comentario{{\scriptsize \(\boxtimes\)} \hyperlink{coro267}{26.7. } }
\end{proposition}

\begin{definition}
    Um espaço \((X,\tau)\) é \emph{localmente conexo} se para cada \( x \in X, \ \exists \B_x \subset \U_x\), base de vizinhanças abertas e conexas. 
\end{definition}

\begin{proposition}
    \textbf{27.6.} \(X\) é localmente conexo \(\ \leftrightharpoons \ \) Componenentes conexas dos abertos são abertas. % e fechadas.  %, \ C_x = \operatorname{int}(C_x) = \overline{C_x} : x \in U \).
\end{proposition}

\demo{
    (\(\Rightarrow\)) Dados \(x \in U\ab\subset X\), \(\forall y \in C_x\) (em \(U\)), \(\exists V\in \B_y: V \subset C_x \ab \subset U\); (\(\Leftarrow\)) Dados \( x \in U\ab \in \U_x, \ C_x\ab \subset U\), em particular, \(C_x \ab \in \U_x\) é conexa. 
}

\begin{corollary}
    \textbf{27.7.} Componentes conexas de localmente conexo são abertas e fechadas. 
\end{corollary}

\begin{proposition}
    \textbf{27.8.} Cada quociente de localmente conexo é localmente conexo. 
\end{proposition}

\demo{
    %Sejam \((X, \tau)\) localmente conexo, 
    Sejam \(q: X \twoheadrightarrow Y\) quociente, \(V\ab \subset Y\) e \(D\subset V\) componente conexa. Dada \(C_x\) (em \(q^{-1}(V)\)) \(: x \in q^{-1}(D)\), temos \(q(C_x)\subset V\) conexo, mas \( q(x) \in D\), logo \(q(C_x) \subset D\) e \(C_x\ab \subset q^{-1}(D)\ab \subset X \) e portanto \(D\ab \subset Y\). 
}

\subsection*{Espaços conexos por caminhos}

\begin{definition}
    Diz-se que \((X,\tau)\) é \emph{conexo por caminhos} se \(\forall a, b \in X, \ \exists f: [0,1] \to X\) contínua (\emph{caminho})  tal que \(f(0) = a\) e \(f(1) = b \). Escrevemos \(f: a \leadsto b\).  
\end{definition}

\begin{proposition}
    \textbf{28.2.} Cada espaço conexo por caminhos é conexo. \comentario{Se \(\exists A \mid B = X\), então basta pegar \(a \in A, \ b\in B\) e \(f: a\leadsto b\), assim \([0,1] = f^{-1}(A) \mid f^{-1} (B )\)\Lightning}
\end{proposition}

\begin{example}
    \(\R^n \setminus E : E \subset \R^n\) é enumeravél é conexo por caminhos. \comentario{Pensa em retas}
\end{example}

\begin{definition}
    Dados \(f: a \leadsto b \) e \(g : b \leadsto c\) em \(X\) definimos \(f *g : a \leadsto c\) tal que:
    \[(f*g) (t) := \begin{cases}
        f(2t), &0\leq t \leq \nicefrac{1}{2} \\ 
        g(2t-1), & \nicefrac{1}{2} \leq t \leq 1
    \end{cases}.\] 
\end{definition}

\vspace{-0.4cm}
\begin{definition}
    Um espaço \((X, \tau)\) é \emph{localmente conexo por caminhos} se cada \( x \in X\) admite uma base de vizinhanças abertas e  conexas por caminhos. 
\end{definition}

\begin{proposition}
    \textbf{28.6.} Se \(X\) é localmente conexo por caminhos, então é conexo por caminhos. \comentario{Seja \(a \in X\), mostra que \(S:= \{x \in X : \exists f:a\leadsto x\}\neq \varnothing\) é aberto e fechado }
\end{proposition}

 %Dado \(b \in S, \ \exists U\ab \in \U_b\) conexa por caminhos, \(\forall c \in U, \ \exists g: b \leadsto c\), e como \(b \in S, \ \exists f: a \leadsto b\), logo \(f*g : a \leadsto c\) e \(U\ab \subset S \ \leftrightharpoons \ S\ab\subset X\).   }

